\section{开源的信仰与实践}

\subsection{为什么坚持开源}

陈天奇是 open source、open science 的坚定信仰者。他认为开源的价值体现在多个层面：

\begin{enumerate}
    \item \textbf{社区交流}：相互分享、认识不同的人、获得代码点赞的正反馈
    \item \textbf{高校落地的路径}：高校本身不容易直接落地技术，开源是分享给大家、共建的最好方式
    \item \textbf{学生培养}：在公司里你通常从大项目的小部件做起，但做开源项目的 PhD 有机会\textbf{从头架构一个系统}------这种训练极其珍贵
    \item \textbf{高校的天然优势}：中立性 + 社区奉献精神
\end{enumerate}

\subsection{构建开源社区的关键}

\begin{importantbox}{成功开源项目的三要素}
\begin{enumerate}
    \item \textbf{正确的技术}：让大家 believe in it。技术本身是基础，没有好的技术，社区不会聚集。
    \item \textbf{持续维护的 passion}：一个软件不是以某个时刻的代码形态决定成功，而是需要被不断维护。Issue 和用户反馈``非常宝贵''------它们是迭代改进的最好输入。
    \item \textbf{Build trust}：如果做了一个东西突然不维护了，大家在下一个项目发布时就会怀疑``他到底是不是需要相信这件事情''。信任需要长期积累，而且一旦失去很难重建。
\end{enumerate}
\end{importantbox}

陈天奇自己会``经常去看'' issue 和用户反馈。他认为开源不是 free 的------需要大量时间看反馈、修 bug、跟社区互动，这些在传统科研评价体系中没有直接回报（不算论文发表）。``所以并不是所有的组几乎应该是很少的组会采取这种路线，但它是一种方式。''

主持人追问：你所有开源项目的 issue 都会一个一个去看吗？陈天奇回答会经常去看，因为``一个软件它不可能是以它某一个时候的代码形态来决定它的成功------它是要被不断维护的''。

\subsection{本章小结}

开源对陈天奇不仅是一种技术传播方式，更是一种价值观：让技术被更多人使用，让每个人都能创造自己的 AI 系统。这也是他做 OctoML 的初心------``为了开源社区能够繁荣''。

